%!TEX root = ../chapter05.tex 
%----------------------------------------------------------------------------------------
%	
%----------------------------------------------------------------------------------------
%----------------------------------------------------------------------------------------
%	 
%----------------------------------------------------------------------------------------

%----------------------------------------------------------------------------------------
%	EXERCICES
%----------------------------------------------------------------------------------------

\newpage 
\loadgeometry{widelayout}  
\pagestyle{scrheadings} % Print headers again  
\KOMAoptions{
	headwidth=textwithmarginpar,
	headsepline=true,
	headsepline= 0.5pt,
	footwidth=textwithmarginpar,
} 
\onehalfspacing
\subsection{Exercices : calculs algébriques et géométrie}
 
\begin{exercise}[Triangles égaux : à l'oral]
	
	Si possible, démontrer pour chaque cas que les triangles sont égaux. Les figures ne sont pas à l'échelle. Utiliser la même couleur pour indiquer les angles et côtés homologues.
	
	\noindent\parbox[position=t]{0.33\textwidth}{
		%\begin{figure}[H]\centering
		\begin{tikzpicture}[scale=0.75]
			\tkzText(0.5,2.8){$(A)$}
			\tkzSetUpPoint[size=3, fill=black]
			\tkzfctset{tan style/.style={->,>=latex, color=red, line width=1.2pt}}   
			\tkzSetUpLine[line width= 1.2pt, add=.5 and .5] 
			\tkzDefPoint(1,0){A} 
			\tkzDefShiftPoint[A](-5:3.3){B}
			\tkzDefTriangle[two angles=90 and 35](A,B)\tkzGetPoint{C}  
			\tkzDrawPolygon[line width=1.2pt](A,B,C)
			%		\tkzLabelPoints[font=\scriptsize](A,B,C)
			\tkzMarkRightAngle[](B,A,C)
			%\tkzMarkAngle[mark=none,-, >=latex, size =0.5](C,B,A)
			%\tkzMarkAngle[mark=none,-, >=latex, size =0.5](A,C,B)
			%\tkzLabelAngle[font=\scriptsize,pos=.8](C,B,A){$\ang{70}$}
			%\tkzLabelAngle[font=\scriptsize,pos=.8](A,C,B){$\ang{55}$}
			%\tkzLabelSegment[above, font=\scriptsize](A,C){$10$} 
			\tkzLabelSegment[below, font=\scriptsize](A,B){\SI{7}{\centi\meter}}
			\tkzLabelSegment[above,sloped,  font=\scriptsize](C,B){\SI{9}{\centi\meter}}
			\begin{scope}[xshift=5cm] 
				\tkzDefPoint(0,0){A}
				\tkzDefShiftPoint[A](+15:2.5){B} 
				\tkzDefTriangle[two angles=90 and 35](A,B)\tkzGetPoint{C}  
				
				\tkzDrawPolygon[line width=1.2pt](A,B,C)
				%	\tkzLabelPoints[font=\scriptsize](A,B,C)
				\tkzMarkRightAngle[](B,A,C)
				%\tkzMarkAngle[mark=none,-, >=latex, size =0.5](C,B,A)
				%\tkzMarkAngle[mark=none,-, >=latex, size =0.5](A,C,B)
				%\tkzLabelAngle[font=\scriptsize,pos=.8](C,B,A){$\ang{70}$}
				%\tkzLabelAngle[font=\scriptsize,pos=.8](A,C,B){$\ang{55}$}
				%\tkzLabelSegment[above, font=\scriptsize](A,C){$10$} 
				\tkzLabelSegment[left, font=\scriptsize](A,C){\SI{7}{\centi\meter}}
				\tkzLabelSegment[above right, font=\scriptsize](C,B){\SI{9}{\centi\meter}}
			\end{scope}
		\end{tikzpicture}   
	}
%	\hfill 
	\parbox[position=t]{0.33\textwidth}{
		%\begin{figure}[H]\centering
		\begin{tikzpicture}[scale=0.75]
			\tkzText(0.5,2){$(B)$}
			\tkzSetUpPoint[size=3, fill=black]
			\tkzfctset{tan style/.style={->,>=latex, color=red, line width=1.2pt}}   
			\tkzSetUpLine[line width= 1.2pt, add=.5 and .5] 
			\tkzDefPoint(0,0){A}
			\tkzDefPoint(3,0){B}
			\tkzDefShiftPoint[A](-26:3.1){B}
			\tkzDefShiftPoint[A](32:4.2){C} 
			\tkzDrawPolygon[line width=1.2pt](A,B,C)
			%		\tkzLabelPoints[font=\scriptsize](A,B,C)
			\tkzMarkAngle[mark=none,-, >=latex, size =0.5](C,B,A)
			\tkzMarkAngle[mark=none,-, >=latex, size =0.5](A,C,B)
			\tkzLabelAngle[font=\scriptsize,pos=.9](C,B,A){$\ang{70}$}
			\tkzLabelAngle[font=\scriptsize,pos=1.0](A,C,B){$\ang{60}$}
			\tkzLabelSegment[above left, font=\scriptsize](A,C){\SI{10}{\centi\meter}}
			%		\tkzDefPointWith[orthogonal,K= -0.65](C,A)
			%		\tkzGetPoint{D}
			%		\tkzDefLine[parallel=through D](A,B) \tkzGetPoint{E}
			%		\tkzDrawLine(D,E)
			%		\tkzDrawLine(A,B)
			%		\tkzLabelLine[pos=1.5,  right](A,B){$(\Delta_1)$}
			%		\tkzLabelLine[pos=1.5,  right](D,E){$(\Delta_2)$}
			%		\tkzMarkRightAngle(D,C,A)
			%		
			%		\tkzDrawSegments[line width= 1.2pt](A,C C,D)
			%		\tkzMarkAngle[mark=none,-, >=latex, size =0.7](B,A,C)
			%		\tkzLabelAngles[pos=1.3,  ](B,A,C){$a$}
			%		\tkzMarkAngle[mark=none,-, >=latex, size =0.7](C,D,E)
			%		\tkzLabelAngles[pos=1.3, ](C,D,E){$b$} 
			\begin{scope}[xshift=4cm] 
				\tkzDefPoint(0,0){A}
				\tkzDefPoint(3,0){B}
				\tkzDefShiftPoint[A](-5:4.1){B}
				\tkzDefShiftPoint[A](48:3.3){C} 
				\tkzDrawPolygon[line width=1.2pt](A,B,C)
				%	\tkzLabelPoints[font=\scriptsize](A,B,C)
				\tkzMarkAngle[mark=none,-, >=latex, size =0.5](B,A,C)
				\tkzMarkAngle[mark=none,-, >=latex, size =0.5](A,C,B)
				\tkzLabelAngle[font=\scriptsize,pos=.9](B,A,C){$\ang{50}$}
				\tkzLabelAngle[font=\scriptsize,pos=.8](A,C,B){$\ang{70}$} 
				\tkzLabelSegment[below, font=\scriptsize](A,B){\SI{10}{\centi\meter}}
			\end{scope}
		\end{tikzpicture}   
	}
%\hfill
\parbox[position=t]{0.33\textwidth}{
		%\begin{figure}[H]\centering
		\begin{tikzpicture}[scale=0.75]
			\tkzText(0.5,2.1){$(C)$}
			\tkzSetUpPoint[size=3, fill=black]
			\tkzfctset{tan style/.style={->,>=latex, color=red, line width=1.2pt}}   
			\tkzSetUpLine[line width= 1.2pt, add=.5 and .5] 
			\tkzDefPoint(0.5,0){A} 
			\tkzDefShiftPoint[A](0:3.7){B}
			\tkzDefTriangle[two angles=-25 and -110](A,B)\tkzGetPoint{C}  
			\tkzDrawPolygon[line width=1.2pt](A,B,C)
			%		\tkzLabelPoints[font=\scriptsize](A,B,C)
			%		\tkzMarkRightAngle[](B,A,C)
			\tkzMarkAngle[mark=none,-, >=latex, size =0.8](C,A,B)
			\tkzMarkAngle[mark=none,-, >=latex, size =0.4](A,B,C)
			\tkzMarkAngle[mark=none,-, >=latex, size =0.6](B,C,A)
			\tkzLabelAngle[font=\scriptsize,pos=1.6](C,A,B){$\ang{20}$}
			\tkzLabelAngle[font=\scriptsize,pos=0.7](A,B,C){$\ang{130}$}
			\tkzLabelAngle[font=\scriptsize,pos=1.2](B,C,A){$\ang{30}$}
			%\tkzLabelSegment[above, font=\scriptsize](A,C){ \SI{13}{\centi\meter}} 
			%		\tkzLabelSegment[above, font=\scriptsize](A,B){\SI{8}{\centi\meter}}
			%		\tkzLabelSegment[above right, font=\scriptsize, sloped](C,B){\SI{9}{\centi\meter}}
			\begin{scope}[xshift=4.5cm, yshift= -0.5cm] 
				\tkzDefPoint(0.5,0){A} 
				\tkzDefShiftPoint[A](-5:3.7){B}
				\tkzDefTriangle[two angles=120 and 30](A,B)\tkzGetPoint{C}  
				\tkzDrawPolygon[line width=1.2pt](A,B,C)
				%		\tkzLabelPoints[font=\scriptsize](A,B,C)
				%		\tkzMarkRightAngle[](B,A,C)
				\tkzMarkAngle[mark=none,-, >=latex, size =0.5](B,A,C)
				\tkzMarkAngle[mark=none,-, >=latex, size =0.8](C,B,A)
				\tkzMarkAngle[mark=none,-, >=latex, size =1.0](A,C,B)
				\tkzLabelAngle[font=\scriptsize,pos=0.8](B,A,C){$\ang{130}$}
				\tkzLabelAngle[font=\scriptsize,pos=1.7](C,B,A){$\ang{20}$}
				\tkzLabelAngle[font=\scriptsize,pos=1.5](A,C,B){$\ang{30}$}
				%	\tkzLabelSegment[left, font=\scriptsize](A,C){\SI{9}{\centi\meter}} 
				%			\tkzLabelSegment[below  , font=\scriptsize](A,B){\SI{8}{\centi\meter}}
				%		\tkzLabelSegment[right, font=\scriptsize](C,B){\SI{8}{\centi\meter}}
			\end{scope}
		\end{tikzpicture}   
	}


\noindent\parbox[position=t]{0.4\textwidth}{
		%\begin{figure}[H]\centering
		\begin{tikzpicture}[scale=0.75]
			\tkzText(0.5,2.8){$(D)$}
			\tkzSetUpPoint[size=3, fill=black]
			\tkzfctset{tan style/.style={->,>=latex, color=red, line width=1.2pt}}   
			\tkzSetUpLine[line width= 1.2pt, add=.5 and .5] 
			\tkzDefPoint(0.5,0){A} 
			\tkzDefShiftPoint[A](0:4.1){B}
			\tkzDefTriangle[two angles=25 and 125](A,B)\tkzGetPoint{C}  
			\tkzDrawPolygon[line width=1.2pt](A,B,C)
			%		\tkzLabelPoints[font=\scriptsize](A,B,C)
			%		\tkzMarkRightAngle[](B,A,C)
			%		\tkzMarkAngle[mark=none,-, >=latex, size =0.6](B,A,C)
			\tkzMarkAngle[mark=none,-, >=latex, size =0.4](C,B,A)
			%		\tkzMarkAngle[mark=none,-, >=latex, size =0.5](A,C,B)
			%		\tkzLabelAngle[font=\scriptsize,pos=1.2](B,A,C){$\ang{20}$}
			\tkzLabelAngle[font=\scriptsize,pos=0.7](C,B,A){$\ang{135}$}
			%		\tkzLabelAngle[font=\scriptsize,pos=.8](A,C,B){$\ang{45}$}
			%		\tkzLabelSegment[above, font=\scriptsize](A,C){ \SI{13}{\centi\meter}} 
			\tkzLabelSegment[below, font=\scriptsize](A,B){\SI{15}{\centi\meter}}
			\tkzLabelSegment[above,sloped, font=\scriptsize](C,B){\SI{8}{\centi\meter}}
			\begin{scope}[xshift=4.5cm, yshift=-1cm] 
				\tkzDefPoint(0.5,0){A} 
				\tkzDefShiftPoint[A](+20:3.1){B}
				\tkzDefTriangle[two angles=35 and 115](A,B)\tkzGetPoint{C}  
				\tkzDrawPolygon[line width=1.2pt](A,B,C)
				%		\tkzLabelPoints[font=\scriptsize](A,B,C)
				%		\tkzMarkRightAngle[](B,A,C)
				\tkzMarkAngle[mark=none,-, >=latex, size =0.7](B,A,C)
				%			\tkzMarkAngle[mark=none,-, >=latex, size =0.5](C,B,A)
				\tkzMarkAngle[mark=none,-, >=latex, size =0.7](A,C,B)
				\tkzLabelAngle[font=\scriptsize,pos=1.5](B,A,C){$\ang{20}$}
				%\tkzLabelAngle[font=\scriptsize,pos=.8](C,B,A){$\ang{35}$}
				\tkzLabelAngle[font=\scriptsize,pos=1.1](A,C,B){$\ang{25}$}
				%		\tkzLabelSegment[left, font=\scriptsize](A,C){$9$} 
				\tkzLabelSegment[below right, font=\scriptsize](A,B){\SI{15}{\centi\meter}}
				\tkzLabelSegment[right, font=\scriptsize](C,B){\SI{8}{\centi\meter}}
			\end{scope}
		\end{tikzpicture}   
	}
%	\hfill 
 \parbox[position=t]{0.30\textwidth}{
		%\begin{figure}[H]\centering
		\begin{tikzpicture}[scale=0.75]
			\tkzText(0.5,2.5){$(E)$}
			\tkzSetUpPoint[size=3, fill=black]
			\tkzfctset{tan style/.style={->,>=latex, color=red, line width=1.2pt}}   
			\tkzSetUpLine[line width= 1.2pt, add=.5 and .5] 
			\tkzDefPoint(0.5,0){A} 
			\tkzDefShiftPoint[A](-20:2.2){B}
			\tkzDefTriangle[two angles=90 and 55](A,B)\tkzGetPoint{C}  
			\tkzDrawPolygon[line width=1.2pt](A,B,C)
			%		\tkzLabelPoints[font=\scriptsize](A,B,C)
			\tkzMarkRightAngle[](B,A,C)
			%		\tkzMarkAngle[mark=none,-, >=latex, size =0.6](B,A,C)
			%		\tkzMarkAngle[mark=none,-, >=latex, size =0.4](C,B,A)
			%		\tkzMarkAngle[mark=none,-, >=latex, size =0.5](A,C,B)
			%		\tkzLabelAngle[font=\scriptsize,pos=1.2](B,A,C){$\ang{20}$}
			%		\tkzLabelAngle[font=\scriptsize,pos=0.6](C,B,A){$\ang{70}$}
			%		\tkzLabelAngle[font=\scriptsize,pos=.8](A,C,B){$\ang{55}$}
			\tkzLabelSegment[left, font=\scriptsize](A,C){ \SI{9}{\centi\meter}} 
			%	\tkzLabelSegment[below, font=\scriptsize](A,B){\SI{10}{\centi\meter}}
			\tkzLabelSegment[above right, font=\scriptsize](C,B){\SI{12}{\centi\meter}}
			\begin{scope}[xshift=3.5cm, yshift=-1cm] 
				\tkzDefPoint(0.5,0){A} 
				\tkzDefShiftPoint[A](30:3.1){B}
				\tkzDefTriangle[two angles=35 and 90](A,B)\tkzGetPoint{C}  
				\tkzDrawPolygon[line width=1.2pt](A,B,C)
				%		\tkzLabelPoints[font=\scriptsize](A,B,C)
				%					\tkzMarkRightAngle[](B,A,C)
				\tkzMarkRightAngle[](C,B,A) 
				%			\tkzMarkAngle[mark=none,-, >=latex, size =0.6](B,A,C)
				%			\tkzMarkAngle[mark=none,-, >=latex, size =0.5](C,B,A)
				%			\tkzMarkAngle[mark=none,-, >=latex, size =0.6](A,C,B)
				%			\tkzLabelAngle[font=\scriptsize,pos=0.9](B,A,C){$\ang{55}$}
				%\tkzLabelAngle[font=\scriptsize,pos=.8](C,B,A){$\ang{35}$}
				%			\tkzLabelAngle[font=\scriptsize,pos=0.9](A,C,B){$\ang{70}$}
				\tkzLabelSegment[above left, font=\scriptsize](A,C){\SI{12}{\centi\meter}} 
				% \tkzLabelSegment[below  , font=\scriptsize](A,B){\SI{10}{\centi\meter}}
				\tkzLabelSegment[above right, font=\scriptsize](C,B){\SI{8}{\centi\meter}}
			\end{scope}
		\end{tikzpicture}   
	}
%	\hfill 
	\parbox[position=t]{0.30\textwidth}{
		%\begin{figure}[H]\centering
		\begin{tikzpicture}[scale=0.75]
			\tkzText(0.5,2.5){$(F)$}
			\tkzSetUpPoint[size=3, fill=black]
			\tkzfctset{tan style/.style={->,>=latex, color=red, line width=1.2pt}}   
			\tkzSetUpLine[line width= 1.2pt, add=.5 and .5] 
			\tkzDefPoint(0.5,0){A} 
			\tkzDefShiftPoint[A](-20:2.5){B}
			\tkzDefTriangle[two angles=80 and 50](A,B)\tkzGetPoint{C}  
			\tkzDrawPolygon[line width=1.2pt](A,B,C)
			%		\tkzLabelPoints[font=\scriptsize](A,B,C)
			%		\tkzMarkRightAngle[](B,A,C)
			%		\tkzMarkAngle[mark=none,-, >=latex, size =0.6](B,A,C)
			\tkzMarkAngle[mark=none,-, >=latex, size =0.4](C,B,A)
			\tkzMarkAngle[mark=none,-, >=latex, size =0.5](A,C,B)
			%		\tkzLabelAngle[font=\scriptsize,pos=1.2](B,A,C){$\ang{20}$}
			\tkzLabelAngle[font=\scriptsize,pos=1.0](C,B,A){$\ang{50}$}
			\tkzLabelAngle[font=\scriptsize,pos=1.0](A,C,B){$\ang{40}$}
			\tkzLabelSegment[left, font=\scriptsize](A,C){ \SI{9}{\centi\meter}} 
			%	\tkzLabelSegment[below, font=\scriptsize](A,B){\SI{10}{\centi\meter}}
			%		\tkzLabelSegment[above right, font=\scriptsize](C,B){\SI{12}{\centi\meter}}
			\begin{scope}[xshift=3.5cm, yshift=+2cm] 
				\tkzDefPoint(0.5,0){A} 
				\tkzDefShiftPoint[A](-60:3.1){B}
				\tkzDefTriangle[two angles=90 and 35](A,B)\tkzGetPoint{C}  
				\tkzDrawPolygon[line width=1.2pt](A,B,C)
				%		\tkzLabelPoints[font=\scriptsize](A,B,C)
				\tkzMarkRightAngle[](B,A,C)
				%			\tkzMarkRightAngle[](C,B,A) 
				%			\tkzMarkAngle[mark=none,-, >=latex, size =0.6](B,A,C)
				%			\tkzMarkAngle[mark=none,-, >=latex, size =0.5](C,B,A)
				%			\tkzMarkAngle[mark=none,-, >=latex, size =0.6](A,C,B)
				%			\tkzLabelAngle[font=\scriptsize,pos=0.9](B,A,C){$\ang{55}$}
				%\tkzLabelAngle[font=\scriptsize,pos=.8](C,B,A){$\ang{35}$}
				%			\tkzLabelAngle[font=\scriptsize,pos=0.9](A,C,B){$\ang{70}$}
				\tkzLabelSegment[above left, font=\scriptsize](A,C){\SI{9}{\centi\meter}} 
				% \tkzLabelSegment[below  , font=\scriptsize](A,B){\SI{10}{\centi\meter}}
				\tkzLabelSegment[above right, font=\scriptsize](C,B){\SI{12}{\centi\meter}}
			\end{scope}
		\end{tikzpicture}   
	} 
\end{exercise}
\vfill  
\begin{exercise}[Révision triangles semblables] \hfill 
	
	
\noindent\parbox{0.4\linewidth}{\centering
\begin{tikzpicture}[scale=1.0] 
	\tkzSetUpPoint[size=3, fill=black]
	\tkzfctset{tan style/.style={->,>=latex, color=red, line width=1.2pt}}   
	\tkzSetUpLine[line width= 1.2pt, add=.5 and .5] 
	\tkzDefPoint(0.0,0){A} 
	\tkzDefShiftPoint[A](-25:3.5){B}
	\tkzDefTriangle[two angles=35 and 70](A,B)\tkzGetPoint{C}  
	\tkzDrawPolygon[line width=1.2pt](A,B,C)
	\tkzLabelPoints[above,font=\scriptsize](C)
	\tkzLabelPoints[right,font=\scriptsize](B)
	\tkzLabelPoints[left,font=\scriptsize](A)
	%	\tkzMarkRightAngle[](B,A,C)
	%			\tkzMarkAngle[mark=||,-, >=latex, size =0.5](B,A,C)
	\tkzMarkAngle[mark=|,-, >=latex, size =0.7](C,B,A)
	\tkzMarkAngle[mark=||,-, >=latex, size =0.5](A,C,B)
	%		\tkzLabelAngle[font=\scriptsize,pos=.8](B,A,C){$\ang{80}$}
	%		\tkzLabelAngle[font=\scriptsize,pos=.8](C,B,A){$\ang{30}$}
	%		\tkzLabelAngle[font=\scriptsize,pos=.8](A,C,B){$\ang{70}$} 
	\tkzLabelSegment[above, sloped , font=\scriptsize](A,C){\SI{14}{\meter}} 
	\tkzLabelSegment[below, sloped , font=\scriptsize](A,B){$x$}
	\tkzLabelSegment[below,sloped ,  font=\scriptsize](C,B){\SI{12}{\meter}}
	\begin{scope}[xshift=4.5cm,yshift=-2cm ] 
		\tkzDefPoint(0.0,0){D} 
		\tkzDefShiftPoint[D](60:2.5){E}
		\tkzDefTriangle[two angles=30 and 75](D,E)\tkzGetPoint{F}  
		\tkzDrawPolygon[line width=1.2pt](D,E,F)
		\tkzLabelPoints[below left,font=\scriptsize](D)
		\tkzLabelPoints[below right,font=\scriptsize](E)
		\tkzLabelPoints[above,font=\scriptsize](F) 
		%			\tkzMarkAngle[mark=none,-, >=latex, size =0.5](E,D,F)
		\tkzMarkAngle[mark=|,-, >=latex, size =0.5](F,E,D)
		\tkzMarkAngle[mark=||,-, >=latex, size =0.5](D,F,E)
		%			\tkzLabelAngle[font=\scriptsize,pos=.8](E,D,F){$\ang{70}$}
		%			\tkzLabelAngle[font=\scriptsize,pos=.8](F,E,D){$\ang{30}$}
		%			\tkzLabelAngle[font=\scriptsize,pos=.8](D,F,E){$\ang{80}$}
		\tkzLabelSegment[below,sloped , font=\scriptsize](D,E){\SI{10}{\centi\meter}}
		\tkzLabelSegment[above,sloped , font=\scriptsize](F,E){\SI{6}{\centi\meter}}
		\tkzLabelSegment[above,sloped , font=\scriptsize](D,F){$y$}
	\end{scope}
\end{tikzpicture}  
}  \parbox{0.55\linewidth}{  
\begin{enumerate}[leftmargin=*, label=\alph*)]
	\item Justifier que les triangles $ACB$ et $FED$ sont semblables.
	\item Écrire les égalités des rapports entre les côtés homologues.
	\item Calculer les longueurs $x$ et $y$.
\end{enumerate}
} 
\end{exercise}
\vfill 
\begin{exercise}\hfill 
	
 \noindent\parbox{0.4\linewidth}{\centering
 	\begin{tikzpicture}[scale=1.0] 
 		\tkzSetUpPoint[size=3, fill=black]
 		\tkzfctset{tan style/.style={->,>=latex, color=red, line width=1.2pt}}   
 		\tkzSetUpLine[line width= 1.2pt, add=.5 and .5] 
 		
 		\tkzDefPoint(0,0){A} 
 		\tkzDefPoint(4,0){B} 
 		\tkzDefPoint(4,3){C} 
 		\tkzDefPoint(0,3){D} 
 		\tkzDrawPolygon[line width=1.2pt](A,B,C,D) 
 		\tkzDefPointBy[projection = onto B--D](C)  \tkzGetPoint{H}
 		\tkzDrawSegments(B,D H,C) 
 		 \tkzDrawSegment[dim={\scriptsize $20$,+10pt,}](D,C)
 		 \tkzDrawSegment[dim={\scriptsize $15$,-10pt,}](B,C)
 		\tkzLabelPoints[above right,font=\scriptsize](C)
 		\tkzLabelPoints[below,font=\scriptsize](B)
 		\tkzLabelPoints[left,font=\scriptsize](A)
 		\tkzLabelPoints[left,font=\scriptsize](D)
 		\tkzLabelPoints[below,font=\scriptsize](H)
 		\tkzLabelSegment[above, sloped , font=\scriptsize](B,H){$c$} 
 		\tkzLabelSegment[above, sloped , font=\scriptsize](H,C){$b$}
 		\tkzLabelSegment[above,sloped ,  font=\scriptsize](D,H){$a$}
 		\tkzMarkRightAngle[size=0.3](C,H,D)
 	\end{tikzpicture}  
 } \parbox{0.55\linewidth}{  
	\begin{enumerate}[leftmargin=*, label=\alph*)]
		\item Calculer la longueur de la diagonale du rectangle $ABCD$.
		\item Justifier que les triangles $ABD$ et $DHC$ sont semblables.
		\item Écrire les égalités des côtés homologues.
		\item Déduire les valeurs de $a$, $b$ et $c$.
	\end{enumerate}
} 
\end{exercise}
\vfill 
%----------------------------------------------------------------------------------------
%	 
%----------------------------------------------------------------------------------------

\begin{exercise}[Quadrature du rectangle] \hfill 
	
	\noindent Sur la figure ci-dessous, le triangle $ABC$ est rectangle en $A$, et $(AH)$ est perpendiculaire à $(BC)$. 
	
	\noindent\parbox[position=t]{0.6 \linewidth}{  
		\begin{enumerate} [leftmargin=*, label=\alph*)]
			\item Démontrer que les triangles $ACH$ et $ABH$ sont semblables. 
			\item Écrire les rapports égaux.
			\item En déduire que $x=\sqrt{ab}$. 
		\end{enumerate}  
	}\hfill
	\noindent\parbox[position=t]{0.35 \linewidth}{  \centering
		\begin{tikzpicture}[scale=0.8]
			\tkzSetUpPoint[size=2, fill=black]
			\tkzfctset{tan style/.style={->,>=latex, color=red, line width=1.2pt}}   
			\tkzSetUpLine[line width= 1.2pt, add=.5 and .5]   
			% 
			\tkzDefPoint(0,0){B} 
			\tkzDefPoint(1.5,0){H}
			\tkzDefPoint(7,0){C}  
			\tkzDefPoint(1.5,3.2){A}  
			\tkzDrawPoints(A,B,C, H) 
			\tkzDrawSegments(A,H H,C C,A A,H A,B B,H)
			\tkzLabelPoints[above,font=\scriptsize](A)
			\tkzLabelPoints[below,font=\scriptsize](B,H,C)
			\tkzLabelSegment[right, font=\scriptsize](A,H){  $x$}
			\tkzLabelSegment[below, font=\scriptsize](H,B){  $a$}
			\tkzLabelSegment[below, font=\scriptsize](H,C){  $b$}
			\tkzMarkRightAngle[ size=0.3](A,H,C) 
			\tkzMarkRightAngle[ size=0.3](B,A,C)	
			%
		\end{tikzpicture}
	}
\end{exercise}  
%----------------------------------------------------------------------------------------
%	 
%----------------------------------------------------------------------------------------

\newpage 
%\begin{exercise}[théorème de l'angle inscrit dans un demi-cercle] \label{chapitre02:exercice:calcullittéraletgeometrie01} 
%	Nous voulons démontrer le théorème suivant :
%\begin{tBox}
%	Si le point $A$ appartient au  cercle de diamètre $[BC]$ alors le triangle $ABC$ est rectangle en $A$.
%	\end{tBox} 
% \noindent\parbox{0.55\linewidth}{
% $A$ est sur le cercle de diamètre $[BC]$ : $OA=OB=OC$. 
%
% On pose $x=\widehat{BAO}$ et $y=\widehat{OAC}$.   
%\begin{enumerate}[leftmargin=*,label=\alph*)]
%	\item Exprimer les angles du triangle $OAB$ à l'aide de $x$.
%	\item En déduire que $\widehat{AOC}=2x$.
%	\item Exprimer $\widehat{AOB}$ à l'aide de $y$.
%	\item Montrer que $x+y=\ang{90}$. 
%\end{enumerate}  
%}\parbox{0.4\linewidth}{\centering
%		\begin{tikzpicture}[scale=1.0]
%			\def \a{3.0};
%			\tkzSetUpPoint[size=3, fill=black]
%			\tkzfctset{tan style/.style={->,>=latex, color=red, line width=1.2pt}} 
%			\tkzDefPoint(0,0){B} 
%			\tkzDefPoint(2*\a,0){C}
%			\tkzDefPoint(\a,0){O}  
%			\tkzDefShiftPoint[O](130:\a){A}
%			\tkzDrawSector[fill = white, line width=1.2pt](O,C)(B)
%			\tkzDrawPolygon[ line width=1.0pt](A,B,C)
%			\tkzDrawSegments[ line width=0.75pt](O,A O,B O,C)
%			\tkzMarkSegments[mark=s||,line     width=0.75pt](O,A O,B O,C)
%			\tkzMarkAngle[mark=|,mkcolor=black, thick, size=0.6, fill=red!20, opacity=0.6](B,A,O)
%			\tkzLabelAngle[pos=1.0,circle](B,A,O){$x$}
%			\tkzMarkAngle[mark=,mkcolor=black, thick, size=0.6, fill=blue!20, opacity=0.6](O,A,C)
%			\tkzLabelAngle[pos=1.1,circle](O,A,C){$y$}
%			%\tkzLabelAngle[pos=1.1,circle](A,C,O){$y$}
%			%\tkzLabelAngle[pos=0.5,circle](O,B,A){$x$}
%			
%			%\tkzLabelAngle[pos=0.9,circle](A,O,B){$2y$}
%			%\tkzLabelAngle[pos=0.5,circle](C,O,A){$2x$}
%			
%			\tkzDrawPoints(A,B,C, O) 
%			\tkzLabelPoints[below,font=\scriptsize](C,O,B) 
%			\tkzLabelPoints[above,font=\scriptsize](A)  
%		\end{tikzpicture}   
%} 
%\end{exercise}
\begin{exercise}[théorème de l'angle au centre] \label{chapitre02:exercice:calcullittéraletgeometrie02}  
	Soit un cercle de centre $O$ passant par les points $A$, $B$ et $C$. Nous voulons démontrer le théorème suivant :
	\begin{tBox}
		Dans un cercle, un angle au centre mesure le double d'un angle inscrit interceptant le même arc.
	\end{tBox}
	
	\noindent\parbox{0.6\linewidth}{
		Pour simplifier on suppose que le centre $O$ est intérieur à l'angle aigu $\widehat{BAC}$.
		Les angles $\widehat{BAC}$ et $\widehat{BOC}$ (intérieur) interceptent le même arc de cercle $ {BC}$.
		
		On pose $x=\widehat{BAO}$ et $y=\widehat{OAC}$.  
		\begin{enumerate}[leftmargin=*,label=\alph*)]
			\item Exprimer les angles du triangle $OAB$ à l'aide de $x$.
			\item Exprimer les angles du triangle $AOC$ à l'aide de $y$. 
			\item Montrer que la mesure de l'angle au centre recherché est égal à $2\widehat{BAC}$. 
		\end{enumerate}  
	}
	\hfill \parbox{0.35\linewidth}{\centering
		\begin{tikzpicture}[scale=1.0]
			\def \a{2.5};
			\tkzSetUpPoint[size=3, fill=black]
			\tkzfctset{tan style/.style={->,>=latex, color=red, line width=1.2pt}} 
			\tkzDefPoint(\a,0){O}  
			\tkzDefShiftPoint[O](-130:\a){B}
			\tkzDefShiftPoint[O](-10:\a){C}  
			\tkzDefShiftPoint[O](120:\a){A}
			\tkzDrawSector[fill = white, line width=1.2pt](O,C)(B)
			%		\tkzDrawPolygon[ line width=1.0pt](A,B,C)
			\tkzDrawArc[line width=1.2pt, dashed](O,B)(C)
			%		\tkzDrawCircle[fill = white, line width=1.2pt](O,A)
			\tkzDrawSegments[ line width=1.0pt](A,B A,C)
			\tkzDrawSegments[ line width=0.75pt](O,A O,B O,C)
			\tkzMarkSegments[mark=s||,line     width=0.75pt](O,A O,B O,C)
			\tkzMarkAngle[mark=none,mkcolor=black, thick, size=0.6, fill=red!20, opacity=0.6](B,A,O)
			\tkzLabelAngle[pos=1.0,circle](B,A,O){$x$}
			\tkzMarkAngle[mark=none,mkcolor=black, thick, size=0.6, fill=blue!20, opacity=0.6](O,A,C)
			\tkzMarkAngle[mark=none,mkcolor=black, thick, size=0.6, fill=blue!20, opacity=0.6](B,O,C)
			\tkzLabelAngle[pos=1.1,circle](B,O,C){$?$}
			\tkzLabelAngle[pos=1.1,circle](O,A,C){$y$}
			%\tkzLabelAngle[pos=1.1,circle](A,C,O){$y$}
			%\tkzLabelAngle[pos=0.5,circle](O,B,A){$x$}
			
			%\tkzLabelAngle[pos=0.9,circle](A,O,B){$2y$}
			%\tkzLabelAngle[pos=0.5,circle](C,O,A){$2x$}
			
			\tkzDrawPoints(A,B,C, O) 
			\tkzLabelPoints[below,font=\scriptsize]( O,B) 
			\tkzLabelPoints[above,font=\scriptsize](A) 
			\tkzLabelPoints[right,font=\scriptsize](C)   
		\end{tikzpicture}  
	} 
\end{exercise}
\begin{theorem}[théorème de l'angle inscrit dans un demi-cercle]
	Si le point $A$ appartient au  cercle de diamètre $[BC]$ alors le triangle $ABC$ est rectangle en $A$.
\end{theorem} 
\noindent
\begin{tikzpicture}[scale=1.0]
	\def \a{2.4};
	\tkzSetUpPoint[size=3, fill=black]
	\tkzfctset{tan style/.style={->,>=latex, color=red, line width=1.2pt}} 
	\tkzDefPoint(0,0){B} 
	\tkzDefPoint(2*\a,0){C}
	\tkzDefPoint(\a,0){O}  
	\tkzDefShiftPoint[O](130:\a){A}
	\tkzDrawSector[fill = white, line width=1.2pt](O,C)(B)
	\tkzDrawPolygon[ line width=1.0pt](A,B,C)
	\tkzDrawSegments[ line width=0.75pt](O,B O,C)
	%		\tkzMarkSegments[mark=s||,line     width=0.75pt](O,A O,B O,C)
	\tkzMarkAngle[<->,>=latex,mkcolor=black, thick, size=0.6, fill=red!20, opacity=0.6](B,O,C)
	\tkzLabelAngle[pos=1.0,circle](B,O,C){\ang{180}}
	\tkzMarkRightAngle[size=0.4](B,A,C)
	%		\tkzMarkAngle[mark=|,mkcolor=black, thick, size=0.6, fill=red!20, opacity=0.6](B,A,O)
	%		\tkzLabelAngle[pos=1.0,circle](B,A,O){$x$}
	%		\tkzMarkAngle[mark=,mkcolor=black, thick, size=0.6, fill=blue!20, opacity=0.6](O,A,C)
	%		\tkzLabelAngle[pos=1.1,circle](O,A,C){$y$}
	%\tkzLabelAngle[pos=1.1,circle](A,C,O){$y$}
	%\tkzLabelAngle[pos=0.5,circle](O,B,A){$x$}
	
	%\tkzLabelAngle[pos=0.9,circle](A,O,B){$2y$}
	%\tkzLabelAngle[pos=0.5,circle](C,O,A){$2x$}
	
	\tkzDrawPoints(A,B,C, O) 
	\tkzLabelPoints[below,font=\scriptsize](C,O,B) 
	\tkzLabelPoints[above,font=\scriptsize](A)  
\end{tikzpicture} \hfill 
\begin{tikzpicture}[scale=1.0]
	\def \a{2.4};
	\tkzSetUpPoint[size=3, fill=black]
	\tkzfctset{tan style/.style={->,>=latex, color=red, line width=1.2pt}} 
	\tkzDefPoint(\a,0){O}  
	\tkzDefShiftPoint[O](-130:\a){B}
	\tkzDefShiftPoint[O](-30:\a){C}  
	\tkzDefShiftPoint[O](120:\a){A}
	\tkzDefShiftPoint[O](10:\a){D}
	\tkzDefShiftPoint[O](65:\a){E}
	\tkzDrawArc[fill = white, line width=1.2pt](O,C)(B)
	%		\tkzDrawPolygon[ line width=1.0pt](A,B,C)
	\tkzDrawArc[line width=1.2pt, dashed](O,B)(C)
	%		\tkzDrawCircle[fill = white, line width=1.2pt](O,A)
	\tkzDrawSegments[ line width=1.0pt](A,B A,C)
	\tkzDrawSegments[ line width=1.0pt](D,B D,C)
	\tkzDrawSegments[ line width=1.0pt](E,B E,C)
	\tkzDrawSegments[ line width=1.0pt,dashed](O,B O,C)
	%	\tkzDrawSegments[ line width=0.75pt](O,A O,B O,C) 
	\tkzMarkAngle[mark=|,mkcolor=black, thick, size=0.6, fill=red!20, opacity=0.6](B,A,C) 
	\tkzMarkAngle[mark=|,mkcolor=black, thick, size=0.6, fill=red!20, opacity=0.6](B,D,C) 
	\tkzMarkAngle[mark=|,mkcolor=black, thick, size=0.6, fill=red!20, opacity=0.6](B,E,C) 
	\tkzMarkAngle[mark=none,mkcolor=black, thick, size=0.6, fill=blue!20, opacity=0.6](B,O,C)
	%	\tkzLabelAngle[pos=1.0,circle](B,A,O){$x$}
	%	\tkzLabelAngle[pos=1.1,circle](B,O,C){$?$}
	%	\tkzLabelAngle[pos=1.1,circle](O,A,C){$y$}
	%\tkzLabelAngle[pos=1.1,circle](A,C,O){$y$}
	%\tkzLabelAngle[pos=0.5,circle](O,B,A){$x$}
	
	%\tkzLabelAngle[pos=0.9,circle](A,O,B){$2y$}
	%\tkzLabelAngle[pos=0.5,circle](C,O,A){$2x$}
	
	\tkzDrawPoints(A,B,C, O) 
	\tkzLabelPoints[below,font=\scriptsize]( O,B) 
	\tkzLabelPoints[above,font=\scriptsize](A) 
	\tkzLabelPoints[right,font=\scriptsize](C)   
\end{tikzpicture}    
\begin{tikzpicture}[scale=1.0]
	\def \a{2.4};
	\tkzSetUpPoint[size=3, fill=black]
	\tkzfctset{tan style/.style={->,>=latex, color=red, line width=1.2pt}} 
	\tkzDefPoint(\a,0){O}  
	\tkzDefShiftPoint[O](-130:\a){B}
	\tkzDefShiftPoint[O](-30:\a){C}  
	\tkzDefShiftPoint[O](-92:\a){A}
	\tkzDefShiftPoint[O](-55:\a){D}
	\tkzDefShiftPoint[O](65:\a){E}
	\tkzDrawArc[fill = white, line width=1.2pt, dashed](O,C)(B)
	%		\tkzDrawPolygon[ line width=1.0pt](A,B,C)
	\tkzDrawArc[line width=1.2pt](O,B)(C)
	%		\tkzDrawCircle[fill = white, line width=1.2pt](O,A)
	\tkzDrawSegments[ line width=1.0pt](A,B A,C)
	\tkzDrawSegments[ line width=1.0pt](D,B D,C)
	%	\tkzDrawSegments[ line width=1.0pt](E,B E,C)
	\tkzDrawSegments[ line width=1.0pt,dashed](O,B O,C)
	%	\tkzDrawSegments[ line width=0.75pt](O,A O,B O,C) 
	\tkzMarkAngle[mark=|,mkcolor=black, thick, size=0.3, fill=red!20, opacity=0.6](C,A,B) 
	\tkzMarkAngle[mark=|,mkcolor=black, thick, size=0.3, fill=red!20, opacity=0.6](C,D,B) 
	%	\tkzMarkAngle[mark=|,mkcolor=black, thick, size=0.6, fill=red!20, opacity=0.6](B,E,C) 
	\tkzMarkAngle[mark=none,mkcolor=black, thick, size=0.6, fill=blue!20, opacity=0.6](C,O,B)
	%	\tkzLabelAngle[pos=1.0,circle](B,A,O){$x$}
	%	\tkzLabelAngle[pos=1.1,circle](B,O,C){$?$}
	%	\tkzLabelAngle[pos=1.1,circle](O,A,C){$y$}
	%\tkzLabelAngle[pos=1.1,circle](A,C,O){$y$}
	%\tkzLabelAngle[pos=0.5,circle](O,B,A){$x$}
	
	%\tkzLabelAngle[pos=0.9,circle](A,O,B){$2y$}
	%\tkzLabelAngle[pos=0.5,circle](C,O,A){$2x$}
	
	\tkzDrawPoints(A,B,C, O) 
	\tkzLabelPoints[below,font=\scriptsize]( O,B) 
	\tkzLabelPoints[below,font=\scriptsize](A) 
	\tkzLabelPoints[right,font=\scriptsize](C)   
\end{tikzpicture}   



\begin{theorem}[Théorème de l'angle inscrit]
	Les angles inscrits interceptant le même arc de cercle ont la même mesure
\end{theorem}
\begin{exercise}\hfill 
	
\noindent\parbox{0.4\linewidth}{
\begin{tikzpicture}[scale=1.0]
	\def \a{2.4};
	\tkzSetUpPoint[size=3, fill=black]
	\tkzfctset{tan style/.style={->,>=latex, color=red, line width=1.2pt}} 
	\tkzDefPoint(\a,0){O} 
	\tkzDrawPoint(O)
	\tkzDrawCircle[R,line width=1.2pt](O,{\a}cm)
	
	
	
	\tkzDefShiftPoint[O](-130:\a){B}
	\tkzDefShiftPoint[O](120:\a){C}  
	\tkzDefShiftPoint[O](26:\a){A}
	\tkzDefShiftPoint[O](-35:\a){D}
	\tkzDrawSegments(A,B C,D)
	\tkzDrawSegments[dashed](C,B A,D)
	\tkzInterLL(A,B)(C,D)\tkzGetPoint{E}
	
	\tkzLabelSegment[above](A,E){\scriptsize$a$}
	\tkzLabelSegment[above](B,E){\scriptsize$b$}
	\tkzLabelSegment[right](C,E){\scriptsize$c$}
	\tkzLabelSegment[right](D,E){\scriptsize$d$}  
	
	\tkzLabelPoints[left,font=\scriptsize](O) 
	\tkzLabelPoints[right,font=\scriptsize](A) 
	\tkzLabelPoints[below left,font=\scriptsize](B) 
	\tkzLabelPoints[above left,font=\scriptsize](C)  
	\tkzLabelPoints[below right,font=\scriptsize](D)   
	\tkzLabelPoints[above,font=\scriptsize](E) 
%	\tkzDefShiftPoint[O](65:\a){E}
%	\tkzDrawArc[fill = white, line width=1.2pt](O,C)(B)
%	%		\tkzDrawPolygon[ line width=1.0pt](A,B,C)
%	\tkzDrawArc[line width=1.2pt, dashed](O,B)(C)
%	%		\tkzDrawCircle[fill = white, line width=1.2pt](O,A)
%	\tkzDrawSegments[ line width=1.0pt](A,B A,C)
%	\tkzDrawSegments[ line width=1.0pt](D,B D,C)
%	\tkzDrawSegments[ line width=1.0pt](E,B E,C)
%	\tkzDrawSegments[ line width=1.0pt,dashed](O,B O,C)
%	%	\tkzDrawSegments[ line width=0.75pt](O,A O,B O,C) 
%	\tkzMarkAngle[mark=|,mkcolor=black, thick, size=0.6, fill=red!20, opacity=0.6](B,A,C) 
%	\tkzMarkAngle[mark=|,mkcolor=black, thick, size=0.6, fill=red!20, opacity=0.6](B,D,C) 
%	\tkzMarkAngle[mark=|,mkcolor=black, thick, size=0.6, fill=red!20, opacity=0.6](B,E,C) 
%	\tkzMarkAngle[mark=none,mkcolor=black, thick, size=0.6, fill=blue!20, opacity=0.6](B,O,C)
%	%	\tkzLabelAngle[pos=1.0,circle](B,A,O){$x$}
%	%	\tkzLabelAngle[pos=1.1,circle](B,O,C){$?$}
%	%	\tkzLabelAngle[pos=1.1,circle](O,A,C){$y$}
%	%\tkzLabelAngle[pos=1.1,circle](A,C,O){$y$}
%	%\tkzLabelAngle[pos=0.5,circle](O,B,A){$x$}
%	
%	%\tkzLabelAngle[pos=0.9,circle](A,O,B){$2y$}
%	%\tkzLabelAngle[pos=0.5,circle](C,O,A){$2x$}
%	
%	\tkzDrawPoints(A,B,C, O)   
\end{tikzpicture}   
}\hfill
\parbox{0.6\linewidth}{
$A$, $B$, $C$ et $D$ sont des points d'un cercle de centre $O$. On suppose que les cordes $[AB]$ et $[CD]$ se coupent en $E$ situé à l'intérieur du cercle.
\begin{enumerate}[label=\alph*), leftmargin=*]
	\item À l'aide du théorème de l'angle inscrit, justifier que $\widehat{EDA}=\widehat{EBD}$.
	\item Montrer que les triangles $EAD$ et $EBC$ sont semblables.
	\item Écrire les égalités des rapports des côtés homologues.
	\item En déduire que $ab=cd$.
\end{enumerate}
}
\end{exercise}

%----------------------------------------------------------------------------------------
%	 
%----------------------------------------------------------------------------------------
 
